\backhead{FURTHER EXERCISES}{khred}
\begin{exlist}
\item A baseline drainage of 120 mm is changed by $+35$ mm by a trigger and by $-12$ mm by a mulch in their own experiments. Find the drainage composed by differences.\exans{143 mm}
\item Engine 1 predicts 1000 held-out runs with RMSE 18 mm for drainage, whose standard deviation is 70 mm. Find the fraction of the variance explained (approximately $1-RMSE^2/s^2$).\exans{0.93}
\item Engine 2 solves the Richards equation with layers of 2 cm over 1.5 m and a time step of 10 min. How many layer-steps does it compute in a season of 150 days?\exans{1.62 million}
\item A composed yield is $Y_0\times1.10\times0.95\times1.07$ with $Y_0=5.2$ Mg\,ha$^{-1}$. Find the yield and the overall ratio.\exans{5.81 Mg\,ha$^{-1}$; 1.118}
\item A layer is re-anchored: its predictions were 6 per cent too high on average. If the shipped prediction of a yield is 6.36 Mg\,ha$^{-1}$, what is the re-anchored value?\exans{6.0 Mg\,ha$^{-1}$}
\end{exlist}
\backhead{ESSAY QUESTIONS}{khred}
{\small\itshape Answer each question in about one page, with equations, a sketch and at least one number from this chapter.}\par\vspace{3pt}
\begin{exlist}
\item Explain the principle of Engine 1. Why can composition by ratios and differences reproduce a complex model, and where does it fail?
\item Describe Engine 2 and the numerical solution of the Richards equation. What controls its accuracy and its speed?
\item Discuss the held-out test of an engine. What does it show, and what can it not show?
\item Explain how the Dashboard reports the results of the engines, and why the messages tab matters.
\item Compare a meta-model with the model it imitates, using SWAFL and Daisy as the example.
\item Discuss the danger of extrapolation with an engine built from an ensemble.
\end{exlist}
